\chapter{Reading the Path}
\label{ch:reading}

A single quadratic program returns one answer. Re-solving on a grid returns a
scatter of them. The homotopy returns the whole curve exactly, and the curve carries
readouts that no single optimum can. This chapter is about those readouts: what to
plot, what it means, what a statistician already knows about it that a portfolio
manager does not, and---the part that requires discipline---which of those meanings
survive the translation of Section~\ref{sec:transfer} and which do not.

It closes with the open problems, because a graduate reader is entitled to know
where the subject stops.

\section{The coefficient profile}

Plot each coordinate of $x(\lambda)$ against $\lambda$. In regression this is the
LARS/LASSO coefficient profile, the picture a statistician reads to see which
variables a model admits and in what order. In finance the same plot, drawn as the
frontier is traced from minimum variance to maximum return, shows exactly how each
position enters, grows, and leaves, with the turning points as entry and exit events:
a membership and turnover map of the efficient set.
\index{coefficient profile}

It is the same plot. Under Theorem~\ref{thm:identity} it \emph{is} the coefficient
profile when $H = X\T X$. The kinks are breakpoints; a coordinate's segment slope is
$\beta_i$ from Lemma~\ref{lem:affine}; a coordinate pinned at a bound is flat with
$\beta_i = 0$.

Two practical readings that the exactness makes trustworthy. First, an investor
operates in a band of volatility rather than over the whole frontier, and the exact
path lets one trace only that band and read the precise weights at every point inside
it, paying only for the turning points the band contains. Second, because segments
are affine in \emph{weight} space, any intermediate portfolio is an exact convex
combination of its two bracketing turning points---no interpolation error, and no
re-solving.

\section{The support size}

Count the coordinates held strictly between their bounds at each point of the path.
In regression this is the number of nonzero coefficients; in finance it is the
number of positions held. Again the same quantity.
\index{support size}

The count falls in a characteristic knee shape: many holdings shed for little added
risk up to a corner, then a long flat tail of one or two concentrated names. An
investor can read off, at any target volatility, how many names the optimum will
hold, or invert the reading and ask what risk a given cardinality buys. That is a
complexity profile of the efficient set, free from the exact path, and the portfolio
literature had no occasion to define it because it had no exact path to read it from.

Here the discipline of Section~\ref{sec:transfer} is required, and the temptation to
be resisted is specific and attractive.

\begin{remark}[Degrees of freedom: what transfers and what does not]
\label{rem:df}
\index{degrees of freedom}
As a \emph{descriptive} complexity measure the count transfers immediately: the
number of positions a frontier portfolio holds strictly between its bounds is a
well-defined effective dimension for that portfolio, computed along the path at no
cost.

The \emph{distributional} statement is another matter. That the support size is an
unbiased estimate of the degrees of freedom of the fit
\cite{zou2007df,tibshirani2012df} is a theorem about the Lagrange-form estimator at a
\emph{fixed penalty} $\lambda$, with the response Gaussian about $X\beta_0$. The
efficient frontier as the Critical Line Algorithm traces it is indexed by the
gross-exposure bound $c$ (equivalently by target return), which is the
\emph{constrained} form; and $c(\lambda)$ depends on the response. So the unbiasedness
does \emph{not} transfer to a frontier point at fixed $c$.

Where it should transfer is the \emph{penalised} portfolio problem,
$\min \tfrac12 w\T\Sig w - \mu\T w + \tau\norm{w}_1$ at fixed leverage penalty
$\tau$---the Brodie et al.\ slice of Section~3.3 rather than the frontier itself:
there the tuning parameter is held fixed as the data vary and the experiment is the
right one. Even there the extension is not immediate, since the general-constraint
case is outside what \cite{tibshirani2012df} establishes. \emph{The leverage axis, not
the frontier, is where to look.}
\end{remark}

The same caution applies to model selection. Because the count is the LASSO's support
size, its knee is recognisably a model-selection elbow, and the devices statistics
uses to choose a penalty by the corner of a complexity-versus-fit curve---the L-curve
criterion \cite{hansen1992}, information criteria on the degrees of freedom---suggest
a principled rule for choosing where to sit on the frontier. \emph{Suggest} is the
right verb: an information criterion needs the degrees-of-freedom statement, which
holds for the penalised parametrisation and not for the bound-indexed frontier.
Turning the elbow into a rule, in the parametrisation where the statistics is valid,
is real work and not bookkeeping.

\section{Selection as the choice of a corner}

Post-selection inference conditions on the event that a particular model was
selected, which for the LASSO is a polyhedron in the response space
\cite{lee2016,tibshirani2016psi}. Under Theorem~\ref{thm:identity} that polyhedron is
precisely the set of data for which a given corner portfolio is optimal, so each
breakpoint of the frontier is a selection event of exactly the form the
selective-inference machinery treats. This much is a statement about the geometry of
the path, which the identity does license.

The caveat is the same one and it is not cosmetic: the conditioning event must be
described in the parametrisation in which the analysis is carried out, so conditioning
on the selected support at fixed $\lambda$ and conditioning on it at a fixed
gross-exposure cap are different events, and the second is the one a portfolio manager
is in. The apparatus should carry over; the polyhedron it conditions on has to be
rewritten first.

We flag this as the most concrete of the transfers the unified view makes available,
and as sketch rather than finished argument.

\section{From one path to a cloud of them}

Because a single path is now cheap---an exact whole curve from one homotopy, and
cheaper still under the operator reading of Chapter~\ref{ch:operator}, where a factor
or kernel form costs $O(nk)$ rather than $O(n^2)$ per step---one can afford to compute
not one but hundreds.

Resample the inputs: draw $(\mu,\Sig)$ from a bootstrap, a Bayesian posterior, or a
rolling estimation window; trace the exact path for each. The ensemble is a
\emph{cloud} of paths rather than a single curve.
\index{resampled frontier}

The exactness is what lets the cloud be read coherently. Each path is affine between
its turning points, so evaluating all of them at a shared abscissa---a target
volatility, or the $\ell_1$ budget of Theorem~\ref{thm:identity}---is \emph{exact}
interpolation, and at each level one obtains a genuine cross-sectional distribution of
weights and values across the draws, not the scatter a grid of separately solved
programs would leave.

Every readout of this chapter then lifts from a point to a distribution. The
coefficient profile acquires a band showing how stable each holding is under
estimation error. The support-size knee becomes a \emph{distribution} of elbows, a
direct picture of how firmly the model-selection operating point is pinned. And the
resampled median path sits systematically \emph{above} the true one in the finance
instance, which is the in-sample optimism a single estimated frontier hides and the
cloud exposes.

The ensemble view of estimation error is itself old---it is Michaud's resampled
efficiency \cite{michaud1999}. What the common scheme adds is that each member is
exact and nearly free, so the cloud is affordable at factor-model scale, and that its
readouts carry the regression meanings above through the identity: an estimation-error
band on a frontier is a confidence statement about a LASSO path read in portfolio
coordinates.

\section{Why exactness is not a numerical nicety}

It is tempting to regard the exact path as a refinement of a grid: nicer, but the
grid would do. It would not, and the reason is specific.

Gridding the parameter and interpolating between grid points cuts \emph{the very
corners} that carry the model-selection information. The breakpoints are where the
active set changes---where a variable enters, where a position is dropped---and they
are precisely what a grid can miss. Coordinate-descent and proximal solvers on a
$\lambda$-grid approximate the path in exactly this way; the homotopy returns the
corners themselves.

There is a numerical detail that makes the difference real rather than notional. The
breakpoints of a large path crowd together, so the ratio test that orders them must
use a tolerance \emph{relative} to the scale of $\lambda$. With it, the path comes back
strictly monotone to floating-point precision; with a fixed absolute tolerance,
distinct breakpoints eventually merge and the parameter steps backwards.
Chapter~\ref{ch:numerics} generalises the lesson.

\section{Open problems}

The family poses its open problems once, which is one of the benefits of seeing it as
one family.

\paragraph{Path length.} The number of breakpoints is near-linear in the problem size
in practice, which is the homotopy's reason for being. Yet Mairal and Yu
\cite{mairal2012} construct LASSO paths whose length is \emph{exponential} in the
dimension, the statistical mirror of the long-known exponential bounds on the critical
regions of a multi-parametric QP \cite{tondel2003}. What separates the benign average
from this worst case, when tracing the whole path beats solving at a single parameter,
and whether the Bland tie-break can be replaced by a rule with a worst-case guarantee,
are open in the shared form. One resolution would serve every instance.

\paragraph{Rank deficiency.} Remark~\ref{rem:rank} set this up. The scheme assumes a
positive-definite reduced Hessian at every visited active set, and for the plain LASSO
that assumption can be removed \cite{tibshirani2013} by tracing the
minimum-$\ell_2$-norm solution, on the strength of an insertion--deletion property of
that selection rather than of the segment structure. No counterpart is known for the
generalized or constrained problem, where the penalty matrix and the constraint system
may be arbitrary but the design may not. Whether the bordered system admits a canonical
selection playing the same role, or whether the arbitrary-design LASSO path is
genuinely a different construction that the common scheme does not contain, is the most
interesting question this material leaves open---and a caution against reading the
scheme as more general than it is.

\paragraph{Constraints and breakpoint counts.} Proposition~\ref{prop:reg} shows the
penalty and constraint parametrisations agree except on flat segments, which are
non-generic. Which constraint systems force them on a positive-measure set, and
whether $p$ inequality rows can \emph{multiply} rather than merely add to the
breakpoint count, are not known.

These are questions about the shared object rather than about portfolios or
regressions, which is the point of having identified it.

\begin{notes}
The readouts of Sections~1--4 follow \cite{schmelzer2026perspective}; the corrections
of Remark~\ref{rem:df} are its own, and are worth reading as a case study in
withdrawing an overreaching claim precisely rather than quietly. Resampled efficiency
is \cite{michaud1999}; the L-curve criterion \cite{hansen1992}; degrees of freedom for
the lasso \cite{zou2007df,tibshirani2012df}; selective inference
\cite{lee2016,tibshirani2016psi}. The exponential path-length construction is
\cite{mairal2012}.
\end{notes}

\begin{exercises}

\begin{exercise}
Show that on a segment the objective value is quadratic in $\lambda$ but the weight
vector is affine, and explain why plotting the frontier in
(variance, return) rather than (volatility, return) coordinates makes each segment an
exact parabolic arc.
\end{exercise}

\begin{exercise}
Prove that any portfolio between two adjacent turning points is an exact convex
combination of them, and give the formula for its variance in terms of the two
endpoints' variances and their cross term.
\end{exercise}

\begin{exercise}
Explain in your own words why $c(\lambda) = \norm{\beta(\lambda)}_1$ being
data-dependent blocks the transfer of a degrees-of-freedom statement, but not the
transfer of a breakpoint-count bound.
\end{exercise}

\begin{exercise}
The maximum-Sharpe portfolio maximises $\mu\T w / \sqrt{w\T\Sig w}$ along the
frontier. Show that on an affine segment $w(t) = (1-t)w_0 + tw_1$ this is an affine
numerator over the square root of a quadratic, and derive the single stationary point
$t^\star$ in closed form. Why does unimodality along the frontier mean that checking
the two segments bracketing the best turning point suffices?
\end{exercise}

\begin{exercise}
Give an example where a grid of $50$ parameter values misses a breakpoint entirely,
and quantify the resulting error in the support size.
\end{exercise}

\begin{exercise}
Explain why a fixed absolute tolerance in the ratio test can make the traced parameter
sequence non-monotone, and design a relative tolerance that prevents it.
\end{exercise}

\begin{exercise}[Implementation]
Trace an exact frontier for $n = 30$ and plot the coefficient profile and the support
size against volatility. Mark the knee. Then re-solve at $50$ grid points with a
general QP solver and overlay: how many corners does the grid miss?
\end{exercise}

\begin{exercise}[Implementation]
Build a cloud: from a fixed population factor model, draw $300$ samples, re-estimate
$(\mu,\Sig)$ from each, trace the exact frontier for every draw, and plot the
percentile envelope of returns at each risk level together with the true frontier.
Confirm that the resampled median sits above the truth, and explain why.
\end{exercise}

\end{exercises}
